%% ===========================================================================
%%  Chapter 2 — ABI / ISA Reference
%% ===========================================================================
\section{ABI / ISA Reference}\label{ch:abi}

All x86-64 cores follow the System~V AMD64 ABI as implemented by macOS, with
two Sakum-specific refinements forced by running under Rosetta on Apple
Silicon.  The AArch64 and RV64 ports follow their respective standard ABIs
(AAPCS64, RV64G calling convention).

\subsection{Registers (x86-64)}

\begin{table}[H]
\centering
\begin{tabularx}{\textwidth}{lX}
\toprule
Group & Registers \\
\midrule
Argument / return & \texttt{rdi, rsi, rdx, rcx, r8, r9} (args 1--6);
  \texttt{rax} (return / syscall number) \\
Caller-saved & \texttt{rax, rcx, rdx, rsi, rdi, r8--r11} \\
Callee-saved & \texttt{rbx, rbp, r12, r13, r14, r15} \\
Stack / frame & \texttt{rsp} (16-byte aligned at every \texttt{call}),
  \texttt{rbp} (frame) \\
Flags & \texttt{rflags} (ZF/SF/CF/OF consumed by branch conditions) \\
\bottomrule
\end{tabularx}
\caption{x86-64 register classes.}
\end{table}

\paragraph{Callee-saved discipline.} Every non-leaf routine pushes
\texttt{rbx, r12--r15} on entry and pops them on exit.  Example from the AI
core:
\begin{lstlisting}[language=SakumAsm,caption={Callee-saved push/pop in walk\_dir (assembly/sakum\_ai.s:196).}]
walk_dir:
    push %rbx
    push %r12
    push %r13
    push %r14
    push %r15
    ...
    pop %r15
    pop %r14
    pop %r13
    pop %r12
    pop %rbx
    ret
\end{lstlisting}

\paragraph{The mandated IR loop counter is \texttt{r12}, not \texttt{rbx}.}
The compiler pipeline's evaluator clobbers \texttt{rbx}
(\texttt{movzx ebx,[rdi]} on the constant-fold path), so an \texttt{rbx}
loop counter corrupts the walk.  \texttt{assembly/README.md} documents this
explicitly.

\subsection{Calling convention}

\begin{enumerate}
  \item Integer / pointer arguments pass in \texttt{rdi, rsi, rdx, rcx, r8,
    r9}; further args spill to the stack.
  \item The stack pointer \texttt{rsp} must be \textbf{16-byte aligned}
    immediately before any \texttt{call} (libc \texttt{printf}/\texttt{exit}
    assume it).  The uniform \texttt{\_main} prologue enforces this:
\begin{lstlisting}[language=SakumAsm,caption={\_main prologue (sakum\_simd.s:10, sakum\_ai.s:74).}]
    push %rbp
    mov  %rsp, %rbp
    and  %rsp, -16        # force 16-byte alignment
    sub  %rsp, 32
\end{lstlisting}
  \item \texttt{rax/rcx/rdx/rsi/rdi} are caller-saved and may be clobbered
    by any helper; save \texttt{rax} around a \texttt{call} whose result
    you need.
  \item Variadic libc calls (\texttt{printf}) consume args in
    \texttt{rdi, rsi, rdx, rcx, r8, r9}; a 2-format-arg call must place the
    first \texttt{\%d} in \texttt{rsi} and the second in \texttt{rdx}
    (a common off-by-one bug placed them in \texttt{rdx, rcx}).
\end{enumerate}

\subsection{Syscalls vs.\ libc}

Sakum uses the \emph{libc} wrappers (prefixed \texttt{\_} on macOS), not
raw \texttt{syscall} instructions, because of the Rosetta constraints below.
The macOS syscall number space is \texttt{0x2000000 + n} (BSD).
Typical bare-metal syscall bases observed in the code:

\begin{lstlisting}[language=SakumAsm,caption={Syscall bases (sakum\_pipe.s:34, sakum\_bramann.s).}]
SYS_READ   = 0x2000000 + 3
SYS_WRITE  = 0x2000000 + 4
SYS_OPEN   = 0x2000000 + 5
SYS_CLOSE  = 0x2000000 + 6
SYS_EXIT   = 0x2000000 + 1
SYS_FORK   = 0x2000000 + 2
SYS_WAIT4  = 0x2000000 + 7
SYS_SOCKET = 0x2000000 + 97
\end{lstlisting}

\subsection{Rosetta-specific gotchas (Apple Silicon, x86\_64 under Rosetta)}
\label{sec:rosetta}

These were discovered empirically and are mandatory reading for any port
work; they are reproduced verbatim from \texttt{assembly/README.md}.

\begin{enumerate}[label=R\arabic*.]
\item \textbf{Never use raw \texttt{syscall}.} Under Rosetta a raw
  \texttt{syscall} raises \texttt{SIGSYS}.  Use the libc wrappers
  (\texttt{\_open\$NOCANCEL}, \texttt{\_opendir\$INODE64},
  \texttt{\_readdir\$INODE64}, \texttt{\_fstat\$INODE64}, \texttt{\_fopen},
  \texttt{\_fprintf}, \texttt{\_printf}, \texttt{\_write},
  \texttt{\_sysctl}, \texttt{\_strncmp}).
\item \textbf{Spurious open returns.} \texttt{\_open}/\texttt{\_open\$NOCANCEL}
  return garbage \texttt{0xA7A5} (42949) for a \emph{missing} file when
  called \emph{deep in the call stack} (after \texttt{opendir}/\texttt{readdir}
  + \texttt{printf}).  \textbf{Always \texttt{\_fstat\$INODE64} the fd;}
  fstat $< 0$ $\Rightarrow$ bogus open $\Rightarrow$ skip.  This filter
  removed all 31 spurious chunk loads in the AI core.
\item \textbf{\texttt{printf}+\texttt{fflush} can hang} in heavy walk loops
  (stdio deadlock).  Avoid debug prints inside \texttt{walk\_dir}; prefer
  the final summary or \texttt{\_write} to fd~2.
\item Keep \texttt{rsp} 16-byte aligned before every \texttt{call}; save
  \texttt{r8}/\texttt{r9} (scratch inside \texttt{ingest\_chunk}, not
  callee-saved).
\end{enumerate}

\subsection{Cross-ISA ports}

\begin{itemize}
  \item \textbf{AArch64} (\texttt{sakum\_tracker\_arm64*.s}): standard AAPCS64;
    libc-based; the NEON variant vectorises the line-scan hot loop
    (\texttt{ld1}/\texttt{dup}/\texttt{cmeq}/\texttt{cnt}).
  \item \textbf{ARMv7-A} (\texttt{sakum\_tracker\_arm32\_sys.s}): libc-free,
    \texttt{svc \#0} with Linux syscall numbers (openat=56, read=63,
    write=64, close=57, exit=93).
  \item \textbf{RV64GC} (\texttt{sakum\_tracker\_riscv64\_sys.s}): libc-free,
    \texttt{ecall} with the same numbers; the RVV variant
    (\texttt{sakum\_tracker\_riscv64\_rvv.s}) vectorises the scan with
    \texttt{vsetvli}/\texttt{vle8.v}/\texttt{vmseq.vx}/\texttt{vfirst.m}.
\end{itemize}
